\subsubsection{Генератор нагрузки}

Loadgen --- сервис создания нагрузки, реализованный на Go. Он подключается к
PostgreSQL или MongoDB, заполняет исходные данные и выполняет операции
нагрузки по выбранному сценарию без использования ORM.

Генератор предоставляет HTTP API: \texttt{/health}, \texttt{/metrics},
\texttt{/clear}, \texttt{/seed} и \texttt{/run}. Express API вызывает эти
методы от имени веб-приложения. Для каждого запуска формируется \texttt{run\_id},
по которому Grafana отделяет метрики одного прогона от другого.

\subsubsection{Метрики loadgen}

Через \texttt{/metrics} loadgen отдает метрики Prometheus. Для каждой метрики
используются метки \texttt{run\_id}, \texttt{model}, \texttt{scenario},
\texttt{operation} и \texttt{stage\_clients}. Это позволяет сравнивать модели,
сценарии и ступени параллелизма на одной панели Grafana.

В эксперименте используются следующие метрики:

\begin{itemize}
  \item \texttt{loadgen\_operation\_duration\_seconds} --- гистограмма времени
        выполнения операции, по ней считаются p50, p95 и p99;
  \item \texttt{loadgen\_operation\_total} --- число выполненных операций с
        разделением на успешные и ошибочные;
  \item \texttt{loadgen\_active\_workers} --- текущее число активных клиентов
        на ступени нагрузки;
  \item \texttt{loadgen\_batch\_size} --- размер пакета для операций записи;
  \item \texttt{loadgen\_operation\_rows\_total} --- количество обработанных
        строк или документов;
  \item \texttt{loadgen\_operation\_bytes\_total} --- примерный объем
        обработанных данных.
\end{itemize}

\subsubsection{Операции нагрузки}

В эксперименте используются пять операций:

\begin{itemize}
  \item \texttt{WRITE\_COMPLEX\_EVENT\_BATCH} -- пакетная запись аналитических
        событий с различным \texttt{payload} и массивом обнаруженных
        объектов;
  \item \texttt{WRITE\_TELEMETRY\_STREAM} -- потоковая запись технической
        телеметрии камер;
  \item \texttt{AGG\_OBJECT\_ACTIVITY\_BY\_AREA} -- подсчет обнаруженных
        объектов по территории, зонам, типам объектов и уровню значимости;
  \item \texttt{AGG\_TELEMETRY\_HEALTH\_WINDOW} -- расчет технического
        состояния камер за заданный интервал времени;
  \item \texttt{READ\_INCIDENT\_TIMELINE} -- чтение списка важных событий
        в зоне вместе с камерами, объектами и телеметрией.
\end{itemize}

Для сравнения используются три профиля нагрузки:

\begin{itemize}
  \item \texttt{write-heavy}: 60\% записи событий, 35\% записи телеметрии,
        5\% операций чтения и аналитики;
  \item \texttt{analytics-heavy}: 30\% записи и 70\% чтения с аналитическими
        агрегациями;
  \item \texttt{balanced}: 70\% записи и 30\% чтения, в том числе по 10\%
        на каждую аналитическую операцию.
\end{itemize}

Профиль наполняется через веса операций. На каждой итерации клиент loadgen
выбирает одну из пяти операций по заданным долям: например, в
\texttt{write-heavy} первые 60\% выборов приходятся на запись событий,
следующие 35\% --- на запись телеметрии, а оставшиеся 5\% --- на чтение.
Данные для операций берутся из подготовленного мира наблюдения: территорий,
зон и камер. Новые события и записи телеметрии создаются уже во время
нагрузочного цикла.

Конкретные параметры операции выбираются из того же набора данных. При записи
событий loadgen выбирает зону, одну из камер этой зоны, тип события,
значимость и содержимое \texttt{payload}. При записи телеметрии выбирается
камера и формируется набор технических показателей. Для агрегаций выбирается
территория, а для чтения важных событий --- зона; во всех запросах чтения
используется заданный интервал времени вокруг момента выполнения операции.

В каждом профиле нагрузка выполняется ступенями по 60 секунд. Число
параллельных клиентов последовательно увеличивается: 1, 5, 10 и 25.
Размер пакета для событий равен 25, для телеметрии -- 50.

\subsubsection{Характер запросов}

Операции нагрузки моделируют типовые операции в системе видеонаблюдения: запись
новых наблюдений, подсчет обнаруженных объектов, проверку технического
состояния камер и поиск важных событий в выбранной зоне.

К примеру, при записи события в базу попадает не только строка о самом событии, но и
связанные с ним камеры, уровень значимости, тип события и параметры
\texttt{payload}. В PostgreSQL JSONB вариативная часть сохраняется одним
полем \texttt{payload}. В нормализованной PostgreSQL то же событие
раскладывается на основную таблицу, таблицу связи с камерами, таблицы деталей
события и таблицу обнаруженных объектов. Поэтому эта операция показывает,
насколько дорого записывать сложное событие.

Характер аналитической нагрузки можно показать на операции
\texttt{AGG\_OBJECT\_ACTIVITY\_BY\_AREA}. Она отвечает на вопрос: какие
объекты фиксировались на выбранной территории за заданный интервал времени.
Результат группируется по зоне, типу объекта и уровню значимости события.
Ниже приведен вариант этого запроса для PostgreSQL JSONB:

\begin{lstlisting}[language=SQL, basicstyle=\ttfamily\small]
SELECT z.zone_code, obj->>'object_type', e.severity,
       count(*), avg((obj->>'confidence')::numeric)
FROM jsonb.event e
JOIN jsonb.zone z ON z.zone_id = e.zone_id
JOIN jsonb.area a ON a.area_id = z.area_id
CROSS JOIN LATERAL jsonb_array_elements(e.payload->'objects') obj
WHERE a.area_code = $1 AND e.occurred_at BETWEEN $2 AND $3
GROUP BY z.zone_code, obj->>'object_type', e.severity;
\end{lstlisting}

Для вложенной MongoDB та же операция выполняется по коллекции
\texttt{events}. Так как территория, зона и обнаруженные объекты уже лежат
внутри документа события, конвейер состоит из фильтрации, раскрытия массива
\texttt{payload.objects} и группировки:

\begin{lstlisting}[basicstyle=\ttfamily\small]
[
  { $match: { "area.area_code": areaCode,
              occurred_at: { $gte: from, $lte: to } } },
  { $unwind: "$payload.objects" },
  { $group: {
      _id: { zone: "$zone.zone_code",
             object_type: "$payload.objects.object_type",
             severity: "$severity" },
      count: { $sum: 1 },
      avg_confidence: { $avg: "$payload.objects.confidence" }
  } }
]
\end{lstlisting}

В этом примере нагрузку создает не только фильтрация по территории и времени,
но и разбор массива \texttt{payload.objects}. В нормализованных моделях такой
же запрос требует дополнительно восстановить связи между событиями, зонами и
территориями; в нормализованной MongoDB для этого используются этапы
\texttt{\$lookup}.

Остальные операции чтения имеют тот же характер: они берут данные за заданный
интервал времени и связывают их с территорией, зоной или камерой. Запрос по
телеметрии считает задержку, потери пакетов, температуру и долю записей со
статусом \texttt{signal\_lost}. Запрос по важным событиям выбирает последние
события высокой значимости в зоне и читает связанный с ними контекст камер.
